%! Unit Launch: Study Design — Unit 1, Lesson 1.0 — Experience & Formalize
%! (Activity · QuickNotes · Application) (Answer Key)
% THIS FILE IS THE BLANK, ANSWERED. Every \answerspace height, every \boxguard count, and
% every \begin{work} block is BYTE-IDENTICAL to experience/main.tex — that identity is what
% keeps the key from drifting off the blank's pagination.
% NO teachernote here — teacher prose goes in the lesson plan.
\documentclass[12pt]{article}
\usepackage{saar-article}
\usepackage{saar-key}   % pulls in -boxes; do NOT also load -boxes
\newcommand{\answerspace}[2]{\par\nopagebreak\noindent\begin{minipage}[t][#1][t]{\linewidth}%
  \color{keyred}\bfseries #2\end{minipage}\par}

\begin{document}

\pageheader{Unit 1, Lesson 1.0}{Experience \& Formalize: Data Detectives --- Answer Key}
% NAMESTRIP: no \namedateperiod here — mirrors the blank.

\vspace{0.06in}

% ── Part 1: the Activity (22 min) ────────────────────────────────────────────
\begin{headlinebox}{frost}
\textbf{Data Detectives.} The Riverbend Athletics Department keeps a record of its varsity
basketball team, and a sports blog just used that record to make a claim about the team. Use
only what you already know, and be ready to show me \textbf{where you see it in the table}.
\end{headlinebox}

\vspace{0.04in}

\begin{scenariobox}[1. The Riverbend Roster]{cerulean}
\begin{center}
\small
\renewcommand{\arraystretch}{1.1}
\begin{tabular}{@{} l c l c c c @{}}
\toprule
\textbf{Player} & \textbf{Jersey \#} & \textbf{Position} & \textbf{Points} &
\textbf{Height (in)} & \textbf{Captain} \\
\midrule
Nia    & 12 & Guard   & 14 & 64 & Yes \\
Ruben  &  7 & Forward &  9 & 71 & No  \\
Sasha  & 23 & Center  &  6 & 75 & No  \\
Marcus &  4 & Guard   & 19 & 68 & Yes \\
Talia  & 17 & Forward & 11 & 70 & Yes \\
Devon  & 21 & Guard   & 13 & 72 & No  \\
\bottomrule
\end{tabular}
\end{center}

\begin{enumerate}[label=\textbf{\alph*.}, leftmargin=1.8em, itemsep=6pt]

  \item How many people does this table describe? \ans{6} \quad Not counting
        ``Player,'' how many things are recorded about each one? \ans{5}

  \item Find the mean number of points.
        \begin{work}
          \bar{x} &= \dfrac{14+9+6+19+11+13}{6} \\
                  &= \dfrac{72}{6} \\
                  &= 12 \text{ points}
        \end{work}

  \item Write one sentence saying what your answer to \textbf{b} describes. Name the group
        and the units.
        \answerspace{1.4cm}{The six players on the Riverbend roster scored a mean of 12
        points each.}

  \item Sasha's position is ``Center.'' Could you find the team's \emph{mean position}?
        Answer, and say what goes wrong if you try.
        \answerspace{1.4cm}{No. Guard, Forward, and Center are names, not amounts, so there
        is nothing to add --- and a mean is built from addition.}

  \item Sort the five recorded columns into two piles: the answer is either an
        \textbf{amount} you could measure or count, or a \textbf{name or label}.

        \begin{center}
        \small
        \renewcommand{\arraystretch}{1.3}
        \begin{tabularx}{0.97\linewidth}{@{} p{2.8cm} >{\raggedright\arraybackslash}p{3.4cm} X @{}}
        \toprule
        \textbf{Column} & \textbf{Amount or label?} & \textbf{How you decided} \\
        \midrule
        Jersey \#   & \ans{label}  & \ans{digits, but they only tell players apart} \\
        Position    & \ans{label}  & \ans{Guard / Forward / Center are names} \\
        Points      & \ans{amount} & \ans{a count of points scored} \\
        Height (in) & \ans{amount} & \ans{a length measured in inches} \\
        Captain     & \ans{label}  & \ans{Yes and No are group names} \\
        \bottomrule
        \end{tabularx}
        \end{center}

  \item A mean is not available for the ``Captain'' column, so summarize it a different way:
        what \textbf{percent} of these six players are captains? \ans{50\%}

\end{enumerate}
\end{scenariobox}

\boxguard[16]
\begin{scenariobox}[2. What the Blog Said]{cerulean}
A sports blog wrote this about the same roster:

\vspace{-2pt}
\begin{center}\vspace{-4pt}
\itshape ``The Riverbend roster's mean jersey number is 14. That is low for a varsity
squad --- this is a young, inexperienced team.''\vspace{-4pt}
\end{center}

\begin{enumerate}[label=\textbf{\alph*.}, leftmargin=1.8em, itemsep=6pt]

  \item Check the blog's arithmetic.
        \begin{work}
          \bar{x} &= \dfrac{12+7+23+4+17+21}{6} \\
                  &= \dfrac{84}{6} \\
                  &= 14
        \end{work}

  \item The arithmetic is correct. \textbf{Is the conclusion?} Look at which pile you put
        ``Jersey \#'' in for part \textbf{e}, and use it to explain what is wrong with the
        blog's claim.
        \answerspace{2.2cm}{No. Jersey number went in the label pile --- the digits only tell
        the players apart, so their mean measures nothing. Nobody wears 14 on average, and
        the roster says nothing at all about age or experience.}

  \item The athletic director has four questions. Mark each one \textbf{O} if \emph{one}
        look-up answers it, or \textbf{M} if you would have to collect answers from
        \emph{many} players.

        \begin{enumerate}[leftmargin=1.8em, itemsep=3pt, topsep=3pt, label=(\roman*)]
          \item What number does Sasha wear? \hfill \ans{O}
          \item How many points do Riverbend players score in a game? \hfill \ans{M}
          \item How tall is Marcus? \hfill \ans{O}
          \item How tall are the players on the Riverbend roster? \hfill \ans{M}
        \end{enumerate}

  \item Pick one question you marked \textbf{O} and rewrite it so that answering it would
        take data from many players.
        \answerspace{1.4cm}{Sample: (iii) becomes ``How tall are the students who try out for
        the Riverbend basketball team?'' Different students give different heights.}

\end{enumerate}
\end{scenariobox}

% ── Part 2: QuickNotes (the debrief fills this, 14 min) — target: half a page ─
\boxguard[14]
\begin{tcolorbox}[enhanced, colback=frost, colframe=cerulean, boxrule=0.8pt, arc=2mm,
    left=3mm, right=3mm, top=1.5mm, bottom=1.5mm, breakable,
    fonttitle=\bfseries\small\color{white}, title={QuickNotes: Naming What You Found},
    attach boxed title to top left={yshift=-2mm, xshift=4mm},
    boxed title style={colback=cerulean, colframe=cerulean, arc=1.5mm}]
\small
\begin{itemize}[leftmargin=1.1em, itemsep=5pt]
  \item \textbf{Data} are numbers or labels plus the \ans{context} that gives them meaning:
        \emph{who}, \emph{what}, and the \ans{units}.
  \item Each row of a data table is one \ans{individual}; each column is one \ans{variable}.
  \item A \ans{quantitative} variable records a measured amount with units --- summarize
        it with a \ans{mean}.
  \item A \ans{categorical} variable records a label or group name --- summarize it with
        counts and \ans{percents}.
  \item \textbf{The test for a column of digits:} not ``does it look like a number?'' but
        \ans{does arithmetic on it mean anything?} Jersey number and ZIP code fail it.
  \item A \ans{statistical} question is one whose answers are expected to \ans{vary}
        from individual to individual, so it takes data to answer.
\end{itemize}
\end{tcolorbox}

% ── Part 3: the Application (worked together, 10 min) ────────────────────────
\boxguard[14]
\begin{notesbox}{Application: The Student Life Survey}
Six Riverbend students filled out the Student Life Survey.
\begin{center}\vspace{-6pt}
\small
\renewcommand{\arraystretch}{1.05}
\begin{tabular}{@{} l c c l c c @{}}
\toprule
\textbf{Student} & \textbf{Grade} & \textbf{Commute (min)} & \textbf{Lunch} &
\textbf{Siblings} & \textbf{Club member} \\
\midrule
Alina   & 10 & 20 & Home       & 2 & Yes \\
Devon   & 11 &  8 & Cafeteria  & 0 & No  \\
Marisol &  9 & 35 & Home       & 3 & Yes \\
Theo    & 12 & 15 & Cafeteria  & 1 & Yes \\
Priya   & 10 &  5 & Off campus & 2 & No  \\
Jamal   & 11 & 25 & Cafeteria  & 4 & Yes \\
\bottomrule
\end{tabular}\vspace{-6pt}
\end{center}

\textbf{Name them.} This table has \ans{6} individuals and \ans{5} variables.

\textbf{Classify.} Commute is \ans{quantitative}; Lunch is \ans{categorical}; Grade is
\ans{categorical}, even though it is written with digits.

\textbf{Compute.} Find the mean commute time for these six students.
\begin{work}
  \bar{x} &= \dfrac{20+8+35+15+5+25}{6} \\
          &= \dfrac{108}{6} \\
          &= 18 \text{ minutes}
\end{work}

\textbf{Interpret.} What does that number say about this group? Name the group and the units.
\answerspace{1.6cm}{The six students surveyed had a mean one-way commute of 18 minutes. It
describes those six students, not all of Riverbend.}

\textbf{What if?} Suppose the survey had recorded \emph{how many minutes} each student waits
in the lunch line instead of \emph{where} they eat. Would that variable's type change, and
which summary becomes available that was not available before? Justify your answer.
\answerspace{1.8cm}{Yes --- minutes in line is a measured amount, so the variable becomes
quantitative and a mean waiting time is available. Place names cannot be added, so no mean
existed before.}
\end{notesbox}

\end{document}
